\section{Introduction}

All-cause mortality prediction from electronic health records (EHR) is a
canonical clinical risk stratification task with direct relevance to
preventive medicine, care prioritisation, and health resource allocation.
Unlike typical classification problems, mortality prediction requires models
to reason over \emph{structured multimodal data}: longitudinal sequences of
disease events with precise temporal ordering, numeric biomarker measurements
with clinical significance thresholds, and demographic context.
Extracting survival-relevant signal from this data is fundamentally a
\emph{structural pattern recognition} problem: identifying which temporal
configurations of ICD codes, and which numeric biomarker profiles, predict
adverse outcomes requires detecting hierarchical, ordered patterns rather
than simple feature co-occurrences.
A disease trajectory is not a bag of tokens: the 8-year interval between a
first metabolic event and a subsequent cardiovascular diagnosis carries causal
information that token proximity alone cannot preserve.
How events are temporally arranged, how age is encoded as a continuous
geometric quantity, and how numeric biomarker values are contextualised
against clinical reference ranges together determine whether a model reads a
patient record as an unordered feature set or as a structured causal
trajectory.
How temporal structure is encoded is the central design question this study
investigates.

Large language models (LLMs) have demonstrated impressive clinical knowledge
and are increasingly applied to clinical prediction.
Two distinct paradigms have emerged: \emph{autoregressive generation},
where a generative EHR transformer produces survival scores directly from
tokenised event sequences~\cite{delphi2024}, and \emph{encoder embedding},
where a frozen encoder maps patient records to dense vectors combined with a
lightweight CoxPH survival head~\cite{llmehrenc,serialehr}.
Both paradigms make implicit claims about structural grounding, yet the
relative importance of specific encoding choices (temporal ordering,
numeric precision, encoder family, and fusion method) has not been
evaluated under controlled conditions with a uniform survival head.

We address this gap with a controlled study on a UK Biobank respiratory
disease subgroup ($n\,=\,124{,}158$), a population where mortality prediction
is especially challenging due to multimorbidity burden and heterogeneous
temporal disease trajectories~\cite{ukbmdrmf}.
Our primary dimension varies \emph{input representation} across five
temporal-encoding settings, from cross-sectional bag-of-features to fully
unified clinical prose, all under a uniform CoxPH survival head.
To verify that findings generalise beyond a single implementation, we include
two ablations: \emph{encoder architecture} (Qwen3-8B vs.\ Bio\_ClinicalBERT)
and \emph{feature fusion method} (concatenation vs.\ element-wise summation).

Our main contributions are:
\begin{itemize}
  \item A five-setting controlled benchmark of temporal encoding strategies
        for LLM-based mortality prediction, the first to isolate input
        serialisation as the sole experimental variable under a uniform
        survival head.

  \item Evidence that \emph{temporal ordering is the critical structural
        property}: decoupled age--event sinusoidal encoding (Setting~4)
        achieves a $+$0.008 C-index advantage (bootstrap 95\,\% CI:
        $+$0.005 to $+$0.010, $p<0.001$) over prose serialisation of
        identical events.

  \item Evidence that LLMs encode numeric biomarker structure from prose:
        unified summary embedding (Setting~5) with \emph{no raw features}
        matches a 39-feature CoxPH baseline, and that autoregressive
        generation faces a calibration gap in acute risk stratification
        (Delphi IBS\,=\,0.185 vs.\ 0.141, near-term AUC\,=\,0.52 at
        9.5\,yr).

  \item Ablation evidence that encoder choice and fusion method are
        secondary to serialisation format: Bio\_ClinicalBERT matches
        Qwen3-8B within 0.001 C-index; concatenation outperforms summation
        by $+$0.028--$+$0.039.
\end{itemize}
